\section{Simplicity with the eigenfunction fixed}
\label{sec:pinning}

The exactification procedure leaves an annulus on which both the metric
and the eigenfunction have their original round form.  This annulus lets
us make the selected eigenvalue simple while preserving the selected
function exactly. Localized trace-free metric variations also occur in
Uhlenbeck's genericity argument \cite[Section~4, Example~2]{Uhlenbeck1976}.
Here the additional requirement that the variation annihilate the
selected gradient preserves the chosen eigenfunction exactly.
We begin with the finite-dimensional point needed for this argument.

\begin{lemma}[A space of symmetric forms]\label{lem:nondegenerate-pencil}
Let $V$ be a finite-dimensional real vector space and let
$\mathcal L\subset\operatorname{Sym}^2(V^*)$ be a linear subspace.
Suppose that, for each $v\ne0$, some $B\in\mathcal L$ satisfies
$B(v,v)\ne0$.  Then $\mathcal L$ contains a nondegenerate form.
\end{lemma}

\begin{proof}
Choose $A\in\mathcal L$ of maximal rank $r$ and suppose that
$\ker A\ne0$.  There is an $r$-dimensional subspace $W$ on which $A$
is nondegenerate.  Fix $0\ne v\in\ker A$ and choose $B\in\mathcal L$
with $B(v,v)\ne0$.  Since $A(v,\cdot)=0$, in a basis of
$W\oplus\R v$ we have
\[
 \det\bigl((A+tB)|_{W\oplus\R v}\bigr)
 =tB(v,v)\det(A|_W)+O(t^2).
\]
This determinant is nonzero for all sufficiently small nonzero $t$.
Thus $A+tB$ has rank at least $r+1$, a contradiction.  The same
argument covers $r=0$, with the empty determinant equal to one.
\end{proof}

\begin{proposition}[Simplicity with the eigenfunction fixed]
\label{prop:pinned-simplicity}
Let $h$ be a smooth metric on $\Sph^3$ and let $u$ be a nonzero real
smooth function satisfying
\[
 -\Delta_hu=\Lambda_nu,\qquad \Lambda_n=n(n+2),\qquad n\ge2.
\]
Suppose that, on an annulus
$\mathcal A=\{a<r<b\}$ with $0<a<b<1$, one has $h=g_*$ and $u=w_n$.
In every $C^\infty$ neighborhood of $h$ there is a smooth metric
$h'$ such that $h'-h$ is supported compactly in $\mathcal A$,
\[
 dV_{h'}=dV_h,\qquad -\Delta_{h'}u=\Lambda_nu,
\]
and $\Lambda_n$ is a simple eigenvalue of $-\Delta_{h'}$.
\end{proposition}

\begin{proof}
All gradients and inner products in this proof use $h$.
On the round annulus, projection of the Euclidean gradient of the
homogeneous harmonic polynomial $w_n$ onto the tangent space of the
sphere gives
\begin{equation}\label{eq:pinning-gradient}
 |\nabla u|^2
 =n^2\bigl(r^{2n-2}-w_n^2\bigr)
 \ge n^2r^{2n-2}(1-r^2)>0.
\end{equation}
Consequently the orthogonal projection $P$ onto
$(\nabla u)^\perp$ is smooth on $\mathcal A$.
Let
\[
 E=\ker(-\Delta_h-\Lambda_n),\qquad
 E_0=\left\{v\in E:\int_{\Sph^3}vu\,dV_h=0\right\}.
\]
These are finite-dimensional real spaces by the compact elliptic
spectral theorem.  If $E_0=0$, there is nothing to prove.

We first show that
\begin{equation}\label{eq:transverse-nonvanishing}
 v\in E_0\setminus\{0\}\quad\Longrightarrow\quad
 P\nabla v\not\equiv0\ \hbox{on }\mathcal A.
\end{equation}
Suppose otherwise.  Then $dv$ is a scalar multiple of $du$ throughout
$\mathcal A$.  Locally write
$x_1+ix_2=re^{i\theta}$, so that $u=r^n\cos(n\theta)$.
The covectors $dr$ and $du$ are independent wherever
$\sin(n\theta)\ne0$.  Such points occur in every annulus under
consideration.  Near one of them, extend $(u,r)$ to a coordinate
system $(u,r,z)$ and restrict to a coordinate product box.  Since
$dv$ annihilates the tangent spaces to the level sets of $u$, its
$r$ and $z$ coordinate derivatives vanish there.  Thus $v=F(u)$ on
this box for a smooth function $F$ of one variable.
The chain rule and the eigenfunction equations give
\begin{equation}\label{eq:pinning-F-equation}
 n^2F''(u)\bigl(r^{2n-2}-u^2\bigr)
 +\Lambda_n\bigl(F(u)-uF'(u)\bigr)=0.
\end{equation}
Differentiating in the coordinate $r$, with $u$ and $z$ fixed, yields
\[
 2n^2(n-1)r^{2n-3}F''(u)=0.
\]
The coefficient is positive because $n\ge2$ and $r>0$.  Therefore
$F''=0$ on the interval in question.  Equation
\eqref{eq:pinning-F-equation} then forces the constant term of $F$
to vanish, so $v=cu$ on this box for some real $c$.
Unique continuation implies $v=cu$ throughout the connected sphere.
Precisely, $v-cu$ solves a second-order elliptic equation with smooth
real positive-definite principal coefficients and smooth lower-order
coefficients; it vanishes on a nonempty open set.  The unique
continuation Theorem of Aronszajn~\cite[pp.~235--236]{Aronszajn1957} applies in these
coordinates and on the connected manifold.  Orthogonality to $u$
now forces $v=0$, proving \eqref{eq:transverse-nonvanishing}.

Let $\mathcal K$ be the real vector space of smooth symmetric
endomorphisms of $T\Sph^3$ supported compactly in $\mathcal A$ and
satisfying
\[
 K\nabla u=0,\qquad \tr K=0.
\]
Each $K\in\mathcal K$ defines a symmetric form on $E_0$ by
\begin{equation}\label{eq:pinning-form}
 M(K)(v,z)=\int_{\Sph^3}\langle K\nabla v,\nabla z\rangle\,dV_h.
\end{equation}
For $v\in E_0\setminus\{0\}$, put $p=P\nabla v$ on $\mathcal A$.
By \eqref{eq:transverse-nonvanishing}, choose a nonnegative smooth
bump function $\chi$ with compact support in $\mathcal A$ such that
$\int\chi|p|^4\,dV_h>0$.  Define, on the annulus and then by zero,
\[
 K=\chi\left(p\otimes p^\flat-\frac{|p|^2}{2}P\right).
\]
Here $(p\otimes p^\flat)X=\langle p,X\rangle p$.
Since $P$ has rank two and $p\perp\nabla u$, this tensor belongs to
$\mathcal K$.  Moreover,
\[
 M(K)(v,v)=\frac12\int_{\mathcal A}\chi|p|^4\,dV_h>0.
\]
The image $M(\mathcal K)$ is a linear subspace of the symmetric
forms on $E_0$.  Lemma~\ref{lem:nondegenerate-pencil} therefore gives
one $K\in\mathcal K$ for which $M(K)$ is nondegenerate on $E_0$.

For real $\tau$, put
\[
 A_\tau=e^{\tau K},\qquad
 h_\tau(X,Y)=h(A_\tau^{-1}X,Y).
\]
The endomorphism $A_\tau$ is positive definite,
$\det A_\tau=e^{\tau\tr K}=1$, and
$A_\tau\nabla u=\nabla u$.  It follows directly that
\[
 dV_{h_\tau}=dV_h,\qquad
 \Delta_{h_\tau}f=\Div_h(A_\tau\nabla f),\qquad
 -\Delta_{h_\tau}u=\Lambda_nu.
\]
Also $h_\tau=h$ off the compact support of $K$, and
$h_\tau\to h$ smoothly as $\tau\to0$.

We claim that $\Lambda_n$ is simple for every sufficiently small
nonzero $\tau$.  If this fails, there are nonzero
$\tau_j\to0$ and real functions $z_j$ such that
\[
 -\Delta_{h_{\tau_j}}z_j=\Lambda_nz_j,\qquad
 \int z_ju\,dV_h=0,\qquad \int z_j^2\,dV_h=1.
\]
The unchanged volume and uniform positivity of $A_{\tau_j}$ give
\[
 \int\langle A_{\tau_j}\nabla z_j,\nabla z_j\rangle\,dV_h
 =\Lambda_n,
\]
so the sequence is bounded in $H^1$.  By Rellich compactness, a
subsequence converges strongly in $L^2$ and weakly in $H^1$ to
$z$.  Passing to the weak eigenfunction equation shows $z\in E_0$;
normalization gives $\|z\|_{L^2(h)}=1$.  For each $v\in E_0$, the
equations for $z_j$ and $v$, tested against each other, imply
\[
 \int\left\langle
 \frac{A_{\tau_j}-\Id}{\tau_j}\nabla z_j,\nabla v
 \right\rangle\,dV_h=0.
\]
The tensor quotient converges uniformly to $K$, so weak $H^1$
convergence gives $M(K)(z,v)=0$ for every $v\in E_0$.
Nondegeneracy of $M(K)$ contradicts $\|z\|_{L^2(h)}=1$.
Choosing any sufficiently small nonzero $\tau$ now completes the proof.
\end{proof}

\section{Persistence and one fixed smooth metric}
\label{sec:persistence}

The lower bound from the wave construction is expressed by a finite
sum of sign-test integrals.  Simplicity makes a strict lower bound
for this functional persist when the metric changes.  We give the
continuity and limiting arguments
explicitly, including the topology in which the limit is taken.

\begin{lemma}[Continuity near a simple eigenvalue]
\label{lem:simple-eigenpair-continuity}
Let $M$ be a closed smooth manifold, let $h_0$ be a smooth Riemannian
metric, and suppose that $\lambda_0>0$ is a simple eigenvalue of
$-\Delta_{h_0}$ with real eigenfunction $u_0$ normalized by
$\|u_0\|_{L^2(h_0)}=1$.  For every open interval $I$ containing
$\lambda_0$, there is a $C^\infty$ neighborhood $\mathcal U$ of
$h_0$ and, for every $h\in\mathcal U$, a simple eigenvalue
$\lambda(h)\in I$ with a normalized real eigenfunction $u(h)$.
One can choose the signs so that
\[
 \lambda(h)\longrightarrow\lambda_0,
 \qquad u(h)\longrightarrow u_0\quad\hbox{in }C^\infty(M)
 \quad\hbox{as }h\longrightarrow h_0.
\]
In particular, if $\mathcal F$ is lower semicontinuous on $C^0(M)$,
invariant under nonzero real rescaling, and $\mathcal F(u_0)>T$,
then $\mathcal U$ can be made small enough that
$\mathcal F(u(h))>T$ for all $h\in\mathcal U$.
\end{lemma}

\begin{proof}
List the eigenvalues with multiplicity as
$0=\mu_0(h)\le\mu_1(h)\le\cdots$, and let $k$ be the index
for which $\mu_k(h_0)=\lambda_0$.  The energy and mass forms are
\[
 Q_h(f)=\int_M |\nabla_hf|_h^2\,dV_h,
 \qquad N_h(f)=\int_Mf^2\,dV_h.
\]
Uniform convergence of the metrics and their inverse tensors implies
that, for some $\varepsilon(h)\to0$,
\[
 (1-\varepsilon)Q_{h_0}\le Q_h\le(1+\varepsilon)Q_{h_0},
 \qquad
 (1-\varepsilon)N_{h_0}\le N_h\le(1+\varepsilon)N_{h_0}.
\]
The variational min--max characterization
\cite[Section~7.2]{CanzaniNotes} therefore gives
$\mu_j(h)\to\mu_j(h_0)$ for every fixed index $j$.
Choose a smaller interval $J$ with closure in $I$ that contains the
selected eigenvalue and no other eigenvalue of $h_0$.  Applying this
convergence at the indices $k-1,k,k+1$ shows that, for $h$ close to
$h_0$, precisely one eigenvalue counted with multiplicity lies in
$J$.  It is $\mu_k(h)$ and it is simple.

For completeness, uniform convergence of the eigenfunctions follows
without differentiating a spectral projection.  Let $h_j\to h_0$
smoothly and let $u_j$ be real eigenfunctions at this index, normalized
in $L^2(h_j)$.  Their eigenvalues stay bounded.  Energy gives an
$H^1(h_0)$ bound.  For every nonnegative integer $s$, the usual global
elliptic estimate for the fixed operator $\Delta_{h_0}$
\cite[Theorem~15.7]{DyatlovNotes} is
\[
 \|f\|_{H^{s+2}}
 \le C_s\bigl(\|\Delta_{h_0}f\|_{H^s}+\|f\|_{L^2}\bigr).
\]
Smooth convergence implies that
$\Delta_{h_j}-\Delta_{h_0}:H^{s+2}\to H^s$ has operator norm tending
to zero.  Absorbing this difference yields the same estimate with
$\Delta_{h_j}$ and a constant independent of large $j$.  The
eigenfunction equation then bounds the $u_j$ in every Sobolev space.
Rellich compactness and Sobolev embedding give subsequential
convergence in every $C^m$ to a smooth eigenfunction of $h_0$ at the
selected eigenvalue.  The limiting $L^2(h_0)$ norm is one because
the volume densities converge uniformly.  Simplicity identifies
this limit as $u_0$ or $-u_0$.
Choosing the sign by
$\int_Mu_ju_0\,dV_{h_0}>0$ excludes the second alternative.  That
integral is nonzero for all sufficiently nearby metrics: otherwise
a sequence on which it vanished would give a normalized limit
orthogonal to $u_0$.  Thus every subsequence has the same limit,
which proves the claimed continuity.
Finally, lower semicontinuity implies
$\liminf\mathcal F(u_j)\ge\mathcal F(u_0)>T$.
The smooth metric topology is metrizable, so the sequential conclusion
is exactly the asserted neighborhood statement.
\end{proof}

We now assemble the preceding constructions into the form needed for
iteration.  Fix a sufficiently small $C^2$ neighborhood
$\mathcal V$ of $g_*$ in the smooth metrics on $\Sph^3$ so that all
uniform comparisons and the initial-profile construction above
apply.  In particular there is one constant $C_{\mathrm{area}}>0$
such that every sign functional furnished by the wave construction
satisfies
\begin{equation}\label{eq:uniform-certificate-area}
 \mathcal F(f)\le C_{\mathrm{area}}\,
 \mathcal H_h^2(\{f=0\}),\qquad h\in\mathcal V,\quad f\in C^0(\Sph^3).
\end{equation}
The constant is independent of the profile, its finite subdivision,
and its frequency; this is Proposition~\ref{prop:sign-functional}.
Each such $\mathcal F$ is lower semicontinuous for uniform convergence
and is invariant under multiplication by any nowhere-zero positive
function, as well as by any nonzero real constant.

\begin{proposition}[One construction step]\label{prop:one-step}
Suppose that $g\in\mathcal V$ agrees with $g_*$ on $r\ge R_0$, where
\[
 \frac15\le R_0<R_1\le\frac3{10}.
\]
Let $\mathcal N$ be any $C^\infty$ neighborhood of $g$, let $A>0$,
and let $N\ge2$ be an integer.  There are an integer $n\ge N$,
a metric $g'\in\mathcal N\cap\mathcal V$ agreeing with $g_*$ on
$r\ge R_1$, and a nonzero real smooth function $u$ such that
$-\Delta_{g'}u=\Lambda_nu$ and $\Lambda_n$ is simple.
Moreover, there is an open neighborhood $\mathcal W$ of $g'$ in
$\mathcal V$ such that every $h\in\mathcal W$ has an exact real
eigenfunction $u_h$ with eigenvalue $\lambda_h$ satisfying
\begin{equation}\label{eq:one-step-open-witness}
 \Lambda_n-1<\lambda_h<\Lambda_n+1,
 \qquad
 \frac{\mathcal H_h^2(\{u_h=0\})}{\sqrt{\lambda_h}}>A.
\end{equation}
This neighborhood is obtained by retaining one fixed sign
functional and one fixed strict numerical threshold.
\end{proposition}

\begin{proof}
Choose a smaller neighborhood $\mathcal N_1$ of $g$ contained in
$\mathcal N\cap\mathcal V$.  Such a neighborhood contains a condition
$\|h-g\|_{C^k}<\epsilon$ for some finite $k$ and $\epsilon>0$.
Choose intermediate radii as in the wave and exactification
constructions, leaving a nonempty annulus between their outermost
modified radius $t_3$ and $R_1$.

Proposition~\ref{prop:spherical-profiles} gives a fixed smooth profile
$\phi$ with
\[
 J:=\int_\Omega L_\phi\,dx
 >4C_{\mathrm{area}}A/c_{\mathrm{wave}}.
\]
Here and below the coordinate measure on the
fixed cube is used; its comparison with $dV_g$ is uniform for
$g\in\mathcal V$.  The wave realization in
Proposition~\ref{prop:wave-realization} supplies, for all sufficiently
large integers $n$, a real function $U$ and a sign functional
$\mathcal F$ with
\[
 \mathcal F(U)\ge c_{\mathrm{wave}}nJ,
\]
where $c_{\mathrm{wave}}>0$ is independent of the chosen profile.
The required finite derivative orders and residual powers can be
specified in advance.  By Proposition~\ref{prop:exactification},
choose these orders so that, after increasing $n$, there is an exact
eigenpair $(\widehat g,u)$ at $\Lambda_n$, with
$\widehat g\in\mathcal N_1$, $u=U/\widetilde v$ for a smooth
function $\widetilde v>0$, and $\widehat g=g_*$, $u=w_n$ for $r\ge t_3<R_1$.
In particular $\mathcal F(u)=\mathcal F(U)$.
All accuracy choices are finite; the profile and these choices are
fixed before $n$ is increased.

Take $n\ge N$ sufficiently large for these conclusions.  Since $n\ge2$,
\begin{equation}\label{eq:one-step-sign-threshold}
 \mathcal F(u)>2C_{\mathrm{area}}A(n+1).
\end{equation}
Apply Proposition~\ref{prop:pinned-simplicity} in an annulus
$t_3<r<t_4<R_1$, with an arbitrarily small smooth change.  It gives
$g'\in\mathcal N_1$ with the same exact eigenfunction $u$ and
simple eigenvalue $\Lambda_n$.  The sign bound is unchanged and
$g'=g_*$ on $r\ge R_1$.

Now freeze $\mathcal F$, including its subdivision and all its
coordinate shifts, and set
$T=C_{\mathrm{area}}A(n+1)$.
Lemma~\ref{lem:simple-eigenpair-continuity} gives an open
neighborhood $\mathcal W\subset\mathcal V$ of $g'$ on which a
normalized real eigenfunction $u_h$ has eigenvalue in
$(\Lambda_n-1,\Lambda_n+1)$ and satisfies
$\mathcal F(u_h)>T$.  The normalization does not affect the sign
functional.  By \eqref{eq:uniform-certificate-area},
\[
 \mathcal H_h^2(\{u_h=0\})>A(n+1).
\]
Finally $\Lambda_n+1=(n+1)^2$, so
$\sqrt{\lambda_h}<n+1$, proving
\eqref{eq:one-step-open-witness}.
\end{proof}

\begin{proof}[Proof of Theorem~\ref{thm:main}]
Let $\mathcal O$ be any prescribed $C^\infty$ neighborhood of the
round metric.  We will retain all the open witness requirements from
Proposition~\ref{prop:one-step} inside $\mathcal O$.

Work in the vector space
$\mathscr T=C^\infty(\Sph^3;\operatorname{Sym}^2T^*\Sph^3)$ of smooth
symmetric covariant tensors.  Fix increasing $C^m$ norms using the
round metric and its connection, and put
\begin{equation}\label{eq:complete-tensor-metric}
 d(T,S)=\sum_{m=0}^\infty 2^{-m-1}
       \min\{1,\|T-S\|_{C^m}\}.
\end{equation}
This is a complete translation-invariant metric inducing the smooth
topology.  Indeed a $d$-Cauchy sequence is Cauchy in every $C^m$
norm, and the compatible $C^m$ limits define one smooth tensor.
Conversely convergence in every such norm implies convergence in
$d$ by splitting the series into a finite sum and a uniformly small
tail.  The same splitting shows that every $d$-neighborhood contains
a neighborhood specified by finitely many seminorm bounds, hence,
since our norms are increasing, by one sufficiently small $C^k$
bound.  This explains why the finite-order closeness supplied by the
construction suffices for the iteration.

The set $\mathcal V$ is open in $\mathscr T$ and consists entirely
of positive-definite tensors.  Set $g_0=g_*$ and choose
$\rho_0>0$ such that the closed $d$-ball
$B_0=\overline B_d(g_0,\rho_0)$ lies in
$\mathcal V\cap\mathcal O$.  Let
\[
 R_j=\frac3{10}-\frac1{10}\,2^{-j},\qquad j\ge0,
\]
so $R_0=1/5$ and $R_j$ increases strictly to $3/10$.
We inductively construct smooth metrics $g_j$, integers $n_j$, open
witness neighborhoods $\mathcal W_j$, and closed balls
$B_j=\overline B_d(g_j,\rho_j)$ with
\begin{align*}
 &g_j=g_*\quad\hbox{on }r\ge R_j,\qquad n_j>n_{j-1},\qquad n_j\ge j+2,\\
 &0<\rho_j<2^{-j},\qquad
 B_j\subset B_d(g_{j-1},\rho_{j-1})\cap\mathcal W_j.
\end{align*}
For $j=1$ the condition involving $n_0$ is omitted.
At step $j$, apply Proposition~\ref{prop:one-step} to $g_{j-1}$,
the pair $(R_{j-1},R_j)$, the open ball
$B_d(g_{j-1},\rho_{j-1})$, the target $A=j$, and a frequency lower
bound larger than $n_{j-1}$ and at least $j+2$.
This supplies $g_j$ in that ball and its open witness neighborhood
$\mathcal W_j$.  Choose $\rho_j>0$ small enough for the displayed
containments and size bound.  This is possible because the center
belongs to both open sets.  At each stage only the constructed center
is required to have a round exterior; the other tensors in $B_j$
need not have one.

The balls are nested and their diameters tend to zero.  Their centers
form a $d$-Cauchy sequence, so completeness gives a smooth limiting
tensor $g_\infty$.  Because each $B_j$ is closed and contains every
later center, $g_\infty\in B_j$ for all $j\ge0$.
In particular $g_\infty\in\mathcal O\cap\mathcal V$ is a smooth
Riemannian metric and belongs to every $\mathcal W_j$.
Each $\mathcal W_j$ retains its own fixed sign functional and
threshold; no uniform spectral gap or common eigenfunction tolerance
is used.

For this one metric, the witness in $\mathcal W_j$ is an exact
nonzero real eigenfunction $u_j$ with
\[
 -\Delta_{g_\infty}u_j=\lambda_ju_j,\qquad
 \Lambda_{n_j}-1<\lambda_j<\Lambda_{n_j}+1,
 \qquad
 \frac{\mathcal H_{g_\infty}^2(\{u_j=0\})}{\sqrt{\lambda_j}}>j.
\]
The integers $n_j$ increase to infinity; the displayed intervals are
pairwise disjoint and ordered, since
$\Lambda_{n+1}-\Lambda_n=2n+3>2$.  Hence
$\lambda_j\to\infty$, while the nodal-area ratios tend to infinity.
The sphere is closed, connected, and three-dimensional, so these are
the required eigenfunctions of one fixed admissible metric.
\end{proof}
